% !TeX root = ../../thesis.tex
\chapter{Conclusions and Future Work}
\label{ch:conclusion}

% ─────────────────────────────────────────────────────────────────────────────
\section{Conclusions}
\label{sec:9_conclusions}

% Paragraph purpose: restate the closed loop concisely, no meta-explanation.
This thesis set out to transform partially scanned indoor environments into complete, structured and dynamically editable \glspl{DRM}. The preceding chapters followed a single artefact through the full modular pipeline of Figure~\ref{fig:1_Research_Aim}: a heterogeneous, partial and out-of-date scan is aligned and updated, decomposed into structural elements and objects with their occlusion maps, completed at both scene and object level, and finally made dynamic.

Chapter~\ref{ch:introduction} framed the central research gap as the absence of a unified framework that jointly addresses alignment, completion and dynamification across heterogeneous and multi-temporal data, and decomposed it into five pipeline objectives together with two cross-cutting objectives, while Chapter~\ref{ch:background} translated the target representation into four concrete \gls{DRM} requirements. Table~\ref{tab:9_synthesis} reports, for each stage, the objective it answers, its input and output, the principal quantitative result on V-Scan and the dominant limitation it carries forward; Table~\ref{tab:9_requirements} maps the same work onto the \gls{DRM} requirements.

\begin{table}[hbtp]
\centering
\caption[Pipeline synthesis]{Synthesis of the pipeline stages: the objective each answers, its input and output, the principal result on V-Scan and the dominant limitation carried forward. Detailed results and experimental conditions are given in the corresponding chapter.}
\label{tab:9_synthesis}
\footnotesize
\renewcommand{\arraystretch}{1.3}
\begin{tabularx}{\textwidth}{@{}>{\bfseries}p{2cm}X@{}}
\toprule
\multicolumn{2}{@{}l}{\textbf{Alignment \& updating (Ch.~\ref{ch:4_Alignment_Updating})}~\cite{vermandere_two-step_2022, vermandere_automatic_2022}} \\
\cmidrule(r){1-1}
Objective & Fuse heterogeneous, multi-temporal scans into one up-to-date frame without full rescanning. \\
Input & Heterogeneous multi-temporal scans \\
Output & Aligned, up-to-date point cloud \\
Results & 0.058\,m average accuracy; 97.5\% empty-scene recovery from two scans. \\
Limitations & Structural geometry aligns more reliably than furnished scenes; scanner changes with low overlap remain hardest. \\
\midrule
\multicolumn{2}{@{}l}{\textbf{Object \& occlusion detection (Ch.~\ref{ch:5_Scene_Segmentation})}~\cite{vermandere_guided_2025, vermandere_synthetic_2026}} \\
\cmidrule(r){1-1}
Objective & Decompose the cloud into structure and objects and measure what is occluded. \\
Input & Aligned point cloud \\
Output & Objects, refined oriented boxes, occlusion grids \\
Results & Primary symmetry axis at $8.01^{\circ}$ error, 94\% finding rate; occluded volume 7.97--22.07\% per class. \\
Limitations & Detection precision of 0.310 propagates false positives downstream, requiring an extra filtering step. \\
\midrule
\multicolumn{2}{@{}l}{\textbf{Scene completion (Ch.~\ref{ch:6_Scene_Completion})}~\cite{vermandere_texture-based_2023,vermandere_semantic_2024}} \\
\cmidrule(r){1-1}
Objective & Reconstruct the static room envelope in geometry and texture. \\
Input & Object-removed cloud, planes, occlusion grid \\
Output & Closed textured mesh of the empty room \\
Results & Chamfer 1.5--1.9\,cm, F-score $>$0.975; PSNR up to 59\,dB. \\
Limitations & Planarity assumption and single-panorama texture limit large or curved occluded regions. \\
\midrule
\multicolumn{2}{@{}l}{\textbf{Object completion (Ch.~\ref{ch:7_Object_Completion})}~\cite{vermandere_geometry_2025, vermandere_guided_2025}} \\
\cmidrule(r){1-1}
Objective & Complete each isolated object using local observations and scene context. \\
Input & Partial object meshes, per-object grids, panorama \\
Output & Fully textured object meshes \\
Results & Voxel method IoU 61.8--72.2\%; image method generalises without per-category tuning. \\
Limitations & Voxel material cap and image-method occlusion-grid decoupling remain. \\
\midrule
\multicolumn{2}{@{}l}{\textbf{Scene dynamification (Ch.~\ref{ch:8_Scene_Dynamification})}~\cite{vermandere_adaptive_2026}} \\
\cmidrule(r){1-1}
Objective & Turn the static reconstruction into an interactive, editable \gls{DRM}. \\
Input & Static scene and object meshes \\
Output & Interactive part-aware \gls{DRM} \\
Results & Selective scaling preserves proportions; openings separated and rendered interactively in Unity. \\
Limitations & Geometry-only part decomposition and planar-wall opening detection remain. \\
\bottomrule
\end{tabularx}
\end{table}

\begin{table}[hbtp]
\centering
\caption[Requirement evaluation]{The four \gls{DRM} requirements of Chapter~\ref{ch:background}, with where each is addressed and the outcome.}
\label{tab:9_requirements}
\footnotesize
\renewcommand{\arraystretch}{1.3}
\begin{tabularx}{\textwidth}{@{}p{4.6cm}p{1.7cm}X@{}}
\toprule
\textbf{Requirement} & \textbf{Addressed in} & \textbf{Outcome} \\
\midrule
Geometric completeness & Ch.~\ref{ch:6_Scene_Completion}, \ref{ch:7_Object_Completion} & Met for structure and objects within the single-room planar scope. \\
Textural completeness & Ch.~\ref{ch:6_Scene_Completion}, \ref{ch:7_Object_Completion} & Partially met: coherent under moderate occlusion; degrades over large unobserved areas. \\
Object-level decomposition & Ch.~\ref{ch:5_Scene_Segmentation} & Met: objects represented as separate entities. \\
Part awareness and scalability & Ch.~\ref{ch:8_Scene_Dynamification} & Met geometrically: convex-part decomposition enables selective scaling; semantics still absent. \\
\bottomrule
\end{tabularx}
\end{table}

% Paragraph purpose: the coupling argument, the scientific core that the tables cannot convey.
The scientific core of these results is the coupling that binds the stages into a single framework rather than a sequence of independent tools, which is the unified treatment Chapter~\ref{ch:introduction} found missing in prior work. Its mechanism is the measured occlusion of Chapter~\ref{ch:5_Scene_Segmentation}: by distinguishing regions that are occupied-but-unobserved from those that are genuinely empty, the occlusion grids let scene completion avoid reconstructing space owned by objects while object completion draws on the surrounding envelope of walls, floors and neighbours. This resolves the mutual dependency that defeats methods treating the two problems in isolation. Validating it quantitatively was itself a methodological obstacle, since the surfaces most relevant to completion are precisely those absent from real scans: the V-Scan dataset of Chapter~\ref{ch:3_Data_Acquisition} overcomes this by pairing realistic partial scans with complete ground truth, and its use as the common benchmark is what lets every stage be checked independently while still composing into one coherent \gls{DRM}.

% Paragraph purpose: cross-cutting findings that recur across stages.
Three findings recur across the stages. First, robustness emerges from combination rather than from any single best method: the alignment vote, the pairing of complementary completion priors and the routing of building elements to different treatments each outperform their strongest individual component. Second, structure is more recoverable than appearance: across alignment, scene completion and dynamification, planar and geometrically regular surfaces are handled reliably, while the residual difficulty concentrates consistently in texture synthesis and in irregular geometry. Third, explicit occlusion reasoning is not only the coupling mechanism but the enabler of modularity itself, since it is what keeps independently developed stages from conflicting when composed.

% Paragraph purpose: honest limitations at the thesis level.
These results are bounded by limitations that the individual chapters identify and that compound at the system level. The scope is restricted to single-room, largely planar indoor environments, so curved structures, multi-room layouts and outdoor scenes fall outside the validated range. Errors propagate forward in the modular design: the moderate detection precision introduces false positives that no later stage corrects, and the planarity and single-panorama assumptions of scene completion limit textural recovery where occlusion is large. Part decomposition for scaling is geometry-only, opening detection depends on planar well-textured walls, and the voxel-resolution trade-off recurs in both occlusion detection and dynamification. Most importantly, the gap towards general deployment in the construction industry remains: the sheer variety of real building contents, materials and capture conditions exceeds what a single furnished-room pipeline, however complete, can currently handle.

% ─────────────────────────────────────────────────────────────────────────────
\section{Future Work}
\label{sec:9_future_work}

% Paragraph purpose: direction 1 — robustness and scope expansion.
The most immediate direction is increasing robustness and expanding the scope beyond the single furnished room. Replacing the current detector with a more recent architecture would reduce the false positives that currently propagate through the pipeline, and extending the planar assumptions of scene completion and opening detection to curved and non-Manhattan structures would broaden the range of buildings the method can handle. Scaling from single rooms to connected multi-room and eventually building-scale environments is the natural next step, and would require the alignment and updating framework to manage inter-room consistency and loop closure rather than a single reference frame.

% Paragraph purpose: direction 2 — fidelity, texture and part awareness.
A second direction is increasing fidelity, particularly in the two areas the cross-cutting analysis identified as residually hard. Texture synthesis under large occlusion would benefit from multi-view and learned scene-level priors rather than a single aligned panorama, and the geometry-only part decomposition used for scaling could be enriched with appearance and learned semantic cues so that part-aware editing reflects function as well as shape. Tightening the link between the image-based object completion and the measured per-object occlusion grid would combine the generality of the image-conditioned generator with the occlusion-awareness of the voxel method, addressing the principal trade-off found in Chapter~\ref{ch:7_Object_Completion}.

% Paragraph purpose: direction 3 — evaluation and the construction-industry gap.
A third direction concerns evaluation and the path towards real deployment. V-Scan establishes controlled ground truth, but closing the gap to the construction industry calls for larger and more varied real-world benchmarks with reliable completion ground truth, captured across the diversity of materials, clutter and sensing conditions encountered in practice. Narrowing the synthetic-to-real domain gap observed in detection and object completion, and validating the full pipeline on such data, is the decisive step from a demonstrated single-room methodology towards the faster, reusable digitisation of existing buildings that motivated this work.

% Paragraph purpose: direction 4 — interaction and downstream use of the DRM.
A fourth direction builds on the interactivity unlocked in Chapter~\ref{ch:8_Scene_Dynamification}. The current dynamification supports repositioning, removal and part-aware scaling, but the editable \gls{DRM} invites richer interaction: physically grounded behaviour, semantic constraints that keep edits plausible, and bidirectional updating in which manual edits and freshly captured scans are reconciled through the alignment and updating framework of Chapter~\ref{ch:4_Alignment_Updating}. Coupling the pipeline to the downstream applications surveyed in Chapter~\ref{ch:Valorisation}, from renovation planning to facility management, would also turn the methodological evaluation used here into task-level evaluation, measuring the \gls{DRM} by the decisions it supports rather than by reconstruction accuracy alone.


%%%%%%%%%%%%%%%%%%%%%%%%%%%%%%%%%%%%%%%%%%%%%%%%%%
% Keep the following \cleardoublepage at the end of this file,
% otherwise \includeonly includes empty pages.
\cleardoublepage